% ========================================================================
%  MinSTEP Objective-C Preprocessor
%  Reference Manual, Version 0.20
%
%  Typeset in TeX to evoke the NeXT Computer, Inc. technical
%  reference manuals of the late 1980s.
% ========================================================================

\documentclass[11pt,twoside]{article}

\usepackage[margin=1in]{geometry}
\usepackage{courier}
\usepackage{url}
\usepackage{tabularx}
\usepackage{longtable}
\usepackage{fancyhdr}

% -----------------------------------------------------------------------
%  Page style
% -----------------------------------------------------------------------

\pagestyle{fancy}
\fancyhf{}
\fancyhead[LE]{\scshape\small MinSTEP Objective-C Preprocessor}
\fancyhead[RO]{\scshape\small Reference Manual}
\fancyfoot[CE]{\small\thepage}
\fancyfoot[CO]{\small\thepage}
\renewcommand{\headrulewidth}{0.4pt}
\renewcommand{\footrulewidth}{0pt}

% -----------------------------------------------------------------------
%  Section formatting (default LaTeX)
% -----------------------------------------------------------------------

% -----------------------------------------------------------------------
%  Commands
% -----------------------------------------------------------------------

\newcommand{\ObjC}{Objective-C}
\newcommand{\objcc}{{\tt objcc}}
\newcommand{\libobjcrt}{{\tt libobjcrt}}
\newcommand{\code}[1]{{\tt #1}}
\newcommand{\meta}[1]{{\itshape #1}}
\newcommand{\ie}{{\itshape i.e.}}

% ========================================================================

\begin{document}

% -----------------------------------------------------------------------
%  Title page
% -----------------------------------------------------------------------

\thispagestyle{empty}
\begin{center}
\vspace*{2in}

{\LARGE\bfseries MinSTEP}\\[0.3in]
{\large\itshape An Open-Source NeXTSTEP 0.8-like Operating Environment}\\[0.5in]
{\LARGE\bfseries Objective-C Preprocessor}\\[0.3in]
{\Large\bfseries Reference Manual}\\[0.8in]
{\large Version 0.20}\\[0.5in]
\vfill
{\large Miguel V.\ Mesquita}\\[0.2in]
{\large July 2026}

\end{center}

\newpage

% -----------------------------------------------------------------------
%  Copyright page
% -----------------------------------------------------------------------

\thispagestyle{empty}
\vspace*{8in}

\begin{center}
\small
Copyright {\copyright} 2026 Miguel V.\ Mesquita.\\
All rights reserved.

\medskip

This document is part of the MinSTEP project and is distributed
under the BSD License.

\medskip

MinSTEP is an open-source project inspired by NeXTSTEP 0.8--2.2.\\
It provides a usermode operating environment with its own kernel
and Mach IPC.

\medskip

\ObjC\ is a trademark of NeXT, Inc.\\
NeXT is a trademark of NeXT, Inc.\\
GCC is a trademark of the Free Software Foundation.

\bigskip

{\bf MinSTEP}\\
Licensed under the BSD License.
\end{center}

\newpage

% -----------------------------------------------------------------------
%  Table of contents
% -----------------------------------------------------------------------

\tableofcontents
\newpage

% ========================================================================
\section{Introduction}
% ========================================================================

This document describes the MinSTEP Objective-C preprocessor, \objcc,
and its companion runtime library, \libobjcrt.  Together they comprise
the MinSTEP Objective-C language toolchain, which translates Objective-C
source files into standard C for compilation by any ANSI-compliant C
compiler.

% ========================================================================
\subsection{Overview}
% ========================================================================

The MinSTEP Objective-C preprocessor is a source-to-source translator.
It reads files written in the Objective-C language and emits equivalent
programs written in pure C.  The generated C code may then be compiled
by a conventional C compiler such as {\tt gcc} to produce machine code.

The translation model is:
\begin{center}
{\tt source.m}~$\longrightarrow$~{\bfseries [objcc]}~$\longrightarrow$~{\tt source.c}~$\longrightarrow$~{\bfseries [cc]}~$\longrightarrow$~{\tt a.out}
\end{center}

This design preserves the full power of the existing C compilation
toolchain.  The preprocessor need only handle the additional constructs
introduced by the Objective-C language; all standard C preprocessing,
compilation, and linking remain unchanged.

The preprocessor is written in standard C (C99) and compiles on any
POSIX-conforming system.  It consists of a single source file of
approximately 1{,}800 lines and requires no external libraries beyond
the standard C library.

% ========================================================================
\subsection{The Objective-C Language}
% ========================================================================

Objective-C is a single-dispatch, dynamically-typed, object-oriented
programming language.  It extends ANSI C with a small number of new
syntactic forms:

\begin{itemize}
  \item Class and category declarations ({\tt @interface},
        {\tt @implementation})
  \item Protocol declarations ({\tt @protocol})
  \item Message expressions ({\tt [receiver selector:argument]})
  \item Instance variable declarations within class definitions
  \item Class methods ({\tt +}) and instance methods ({\tt -})
  \item Dynamic type queries ({\tt @selector}, {\tt @encode})
  \item Automatic property synthesis ({\tt @property},
        {\tt @synthesize})
\end{itemize}

All other constructs are standard C.  The preprocessor's task is
therefore modest: it must recognize the Objective-C extensions, emit C
code of equivalent semantics, and pass all remaining C source text
through unchanged.

% ========================================================================
\subsection{Notation}
% ========================================================================

In this manual, monospaced text ({\tt like this}) indicates language
constructs, filenames, or command-line syntax.  Nonterminal names in
the grammatical descriptions are shown in {\itshape italics}.

% ========================================================================
\newpage
\section{Using the Preprocessor}
% ========================================================================

% ========================================================================
\subsection{Invocation}
% ========================================================================

The preprocessor is invoked as:
\begin{quote}
{\tt objcc} {\meta{options}} {\meta{file.m}} [{\meta{file2.m}~\ldots}]
\end{quote}

The following options are recognized:

\medskip
\noindent
\begin{tabularx}{\textwidth}{lX}
{\tt -E} & Preprocess only.  The translated C source is written to the
           standard output, or to the file named by {\tt -o}. \\[4pt]
{\tt -c} & Compile to an object file.  The preprocessor output is
           passed to the C compiler with the {\tt -c} flag.  This is
           the default mode when no output file is specified. \\[4pt]
{\tt -S} & Compile to assembly language.  The preprocessor output is
           passed to the C compiler with the {\tt -S} flag. \\[4pt]
{\tt -o} {\meta{file}}
         & Write output to {\meta{file}}.  The output mode is inferred
           from the extension of {\meta{file}}:
           \begin{itemize}
             \setlength{\itemsep}{0pt}
             \setlength{\parskip}{0pt}
             \item {\tt .c} $\rightarrow$ preprocess only ({\tt -E})
             \item {\tt .o} $\rightarrow$ compile to object ({\tt -c})
             \item {\tt .s} $\rightarrow$ compile to assembly ({\tt -S})
             \item otherwise $\rightarrow$ compile and link
           \end{itemize} \\[4pt]
{\tt -I} {\meta{dir}}
         & Add {\meta{dir}} to the include search path.  This path is
           used by both the preprocessor's {\tt \#import} directive and
           the underlying C compiler. \\[4pt]
{\tt -D} {\meta{name=val}}
         & Define a C preprocessor macro.  Passed through to the C
           compiler unchanged. \\[4pt]
{\tt -L} {\meta{dir}}
         & Add {\meta{dir}} to the library search path.  Passed to the
           linker. \\[4pt]
{\tt -l} {\meta{lib}}
         & Link against library {\meta{lib}}.  Passed to the linker.
           \\[4pt]
{\tt -W} {\meta{flag}}
         & Pass a warning flag to the C compiler. \\[4pt]
{\tt -std=} {\meta{std}}
         & Pass a language standard flag to the C compiler. \\[4pt]
{\tt -g}  & Include debugging symbols. \\[4pt]
{\tt -O} {\meta{n}}
         & Set optimization level {\meta{n}}. \\[4pt]
{\tt -f} {\meta{flag}}
         & Pass an arbitrary flag to the C compiler. \\[4pt]
{\tt -h}, {\tt --help}
         & Display a usage summary and exit. \\[4pt]
{\tt -v}, {\tt --version}
         & Display the version number and exit. \\
\end{tabularx}

% ========================================================================
\subsection{Environment}
% ========================================================================

The C compiler used for the {\tt -c}, {\tt -S}, and link modes is
determined by the following search:

\begin{enumerate}
  \item The {\tt OBJCC\_CC} environment variable, if set.
  \item The {\tt CC} environment variable, if set.
  \item The string {\tt ``gcc''}, as a final default.
\end{enumerate}

The include search path may be extended by multiple {\tt -I} flags.
These paths are searched in order after the current directory when
resolving {\tt \#import} and {\tt \#include} directives with quoted
filenames.

% ========================================================================
\subsection{Examples}
% ========================================================================

Preprocess a single file and examine the output:
\begin{quote}
{\tt \$ objcc -E MyClass.m -o MyClass.c}
\end{quote}

Compile an Objective-C source file to an object file:
\begin{quote}
{\tt \$ objcc -c MyClass.m -o MyClass.o}
\end{quote}

Compile and link a complete program:
\begin{quote}
{\tt \$ objcc -o myapp main.m MyClass.m -lobjcrt}
\end{quote}

Compile with a custom include path:
\begin{quote}
{\tt \$ objcc -I /usr/local/include/objc -c main.m}
\end{quote}

% ========================================================================
\newpage
\section{Objective-C Source Language Constructs}
% ========================================================================

This section describes each Objective-C construct as it is translated
by the preprocessor.  For each construct, the input syntax and the
generated C output are given.

% ========================================================================
\subsection{Class Declarations}
% ========================================================================

An Objective-C class is declared with the {\tt @interface} directive.
The full syntax is:
\begin{quote}
{\tt @interface} {\meta{ClassName}} [: {\meta{SuperclassName}} [{\meta{ProtocolName,~\ldots}}]]\\
\{\\
\quad {\meta{ivar-declarations}}\\
\}\\
{\meta{method-declarations}}\\
{\tt @end}
\end{quote}

The preprocessor emits a C {\tt struct} whose first member is an {\tt isa}
pointer to the class metaclass and whose second member is a pointer to the
superclass struct.  Instance variables follow.

{\bf Input:}
\begin{quote}
\begin{verbatim}
@interface Rectangle : Shape
{
    int width;
    int height;
    @protected
    id fillColor;
}
- (int)area;
@end
\end{verbatim}
\end{quote}

{\bf Output:}
\begin{quote}
\begin{verbatim}
struct _Rectangle_class_t;
struct Rectangle {
    struct _Rectangle_class_t *isa;
    struct Shape *_super;
    int width;
    int height;
    /* @protected */
    id fillColor;
};
typedef struct Rectangle Rectangle;
- (int)Rectangle_inst_area_(id self, SEL _cmd);
\end{verbatim}
\end{quote}

The {\tt isa} and {\tt \_super} pointers are always present, even for
root classes.  The visibility keywords {\tt @public}, {\tt @protected},
{\tt @private}, and {\tt @package} are converted to C comments;
visibility is not enforced at the C level but is recorded for the
runtime.

% ========================================================================
\subsection{Forward Class Declarations}
% ========================================================================

The {\tt @class} directive introduces a forward declaration:
\begin{quote}
{\tt @class} {\meta{ClassName}};
\end{quote}

This is translated to a C forward struct declaration:
\begin{quote}
{\tt struct} {\meta{ClassName}}; \hfill {\tt /* forward class declaration */}
\end{quote}

The class name is also recorded internally so that subsequent
{\tt @interface} declarations for the same class are not duplicated.

% ========================================================================
\subsection{Implementations}
% ========================================================================

The {\tt @implementation} directive marks the beginning of a class
implementation:
\begin{quote}
{\tt @implementation} {\meta{ClassName}} [({\meta{CategoryName}})]\\
{\meta{method-definitions}}\\
{\tt @end}
\end{quote}

When a category name is present in parentheses, the preprocessor enters
category mode and records both the class and category names.  Method
definitions within the implementation block are translated as described
in Section~5.4.

{\bf Input:}
\begin{quote}
\begin{verbatim}
@implementation MyClass
\end{verbatim}
\end{quote}

{\bf Output:}
\begin{quote}
\begin{verbatim}
/* @implementation MyClass */
\end{verbatim}
\end{quote}

% ========================================================================
\subsection{Protocols}
% ========================================================================

A protocol declares an interface that classes may adopt:
\begin{quote}
{\tt @protocol} {\meta{ProtocolName}} [{\tt <}{\meta{SuperProtocol}}{\tt >}]\\
{\meta{method-declarations}}\\
{\tt @end}
\end{quote}

The preprocessor emits a static {\tt Protocol} pointer:
\begin{quote}
{\tt static Protocol}~{\tt *}{\tt \_proto\_}{\meta{ProtocolName}}~{\tt = NULL;}
\end{quote}

Protocol inheritance is parsed but not expanded; the superclass
relationship is recorded for the runtime to resolve.

% ========================================================================
\subsection{Categories}
% ========================================================================

A category adds methods to an existing class.  The syntax is:
\begin{quote}
\begin{tabular}{l}
{\tt @interface} {\meta{ClassName}} ({\meta{CategoryName}})\\
{\meta{method-declarations}}\\
{\tt @end}
\end{tabular}
\end{quote}
\begin{quote}
\begin{tabular}{l}
{\tt @implementation} {\meta{ClassName}} ({\meta{CategoryName}})\\
{\meta{method-definitions}}\\
{\tt @end}
\end{tabular}
\end{quote}

The preprocessor tracks the category name separately from the class
name.  Method definitions in a category implementation are translated
with the category class name, allowing the runtime to associate the
methods with the correct class.

% ========================================================================
\subsection{Property Declarations}
% ========================================================================

The {\tt @property} directive declares a property with optional attribute
specifiers:
\begin{quote}
{\tt @property} ({\meta{attributes}}) {\meta{Type}} {\meta{name}};
\end{quote}

The following attributes are recognized:

\medskip
\noindent
\begin{tabularx}{\textwidth}{lX}
{\tt readonly}     & The property has no setter. \\
{\tt readwrite}    & The property has both getter and setter (default). \\
{\tt nonatomic}    & The setter is not atomic. \\
{\tt retain}       & The setter retains the assigned value. \\
{\tt copy}         & The setter copies the assigned value. \\
{\tt assign}       & The setter assigns directly (default). \\
{\tt getter=} {\meta{name}}
                   & Use {\meta{name}} for the getter method. \\
{\tt setter=} {\meta{name}}{\tt :}
                   & Use {\meta{name}}{\tt :} for the setter method. \\
\end{tabularx}

The preprocessor emits declarations for the getter and setter methods:

{\bf Input:}
\begin{quote}
\begin{verbatim}
@property (nonatomic, retain) NSString *name;
\end{verbatim}
\end{quote}

{\bf Output:}
\begin{quote}
\begin{verbatim}
/* @property (nonatomic, retain) */
- (NSString *)name;
- (void)setName:(NSString *)value;
\end{verbatim}
\end{quote}

For readonly properties, only the getter is emitted.

{\bf Input:}
\begin{quote}
\begin{verbatim}
@property (readonly) int count;
\end{verbatim}
\end{quote}

{\bf Output:}
\begin{quote}
\begin{verbatim}
/* @property (readonly) */
- (int)count;
\end{verbatim}
\end{quote}

% ========================================================================
\subsection{Property Synthesis}
% ========================================================================

The {\tt @synthesize} directive requests automatic generation of getter
and setter implementations:
\begin{quote}
{\tt @synthesize} {\meta{name}} = {\meta{\_name}};
\end{quote}

The preprocessor looks up the property declared in the current class
and emits the corresponding method bodies:
\begin{quote}
\begin{verbatim}
- (NSString *)name {
    return self->_name;
}
- (void)setName:(NSString *)value {
    self->_name = value;
}
\end{verbatim}
\end{quote}

If the property is readonly, only the getter is generated.  The ivar
name may be customized with the {\tt =} syntax; if omitted, the property
name prefixed with an underscore is used.

The {\tt @dynamic} directive suppresses synthesis and emits a comment:
\begin{quote}
{\tt @dynamic} {\meta{name}};
\end{quote}
which produces:
\begin{quote}
{\tt /* @dynamic} {\meta{name}} {\tt - provided at runtime */}
\end{quote}

% ========================================================================
\subsection{Method Declarations}
% ========================================================================

Methods are declared with a leading {\tt +} (class method) or {\tt -}
(instance method), followed by a type specification and selector:
\begin{quote}
{\tt +} ({\meta{ReturnType}}) {\meta{methodName}}:({\meta{ArgType}}) {\meta{argument}};\\
{\tt -} ({\meta{ReturnType}}) {\meta{methodName}}:({\meta{ArgType}}) {\meta{argument}};
\end{quote}

Multi-keyword selectors are supported:
\begin{quote}
{\tt -} ({\tt id}) {\tt initWithFoo:}({\tt int}) {\tt foo} {\tt bar:}({\tt id}) {\tt bar};
\end{quote}

Each keyword-argument pair introduces a new parameter.  If an argument
type is omitted, {\tt id} is assumed.

The preprocessor translates method names to a flat C naming convention.
The class name, a discriminator ({\tt inst} or {\tt cls}), the method
name, and each argument label are joined with underscores.  A trailing
underscore terminates the name.

The naming rules are:
\begin{quote}
\begin{tabular}{ll}
Instance method: & {\tt ClassName\_inst\_methodName\_}({\tt id self}, {\tt SEL \_cmd}, \\
                 & \quad {\meta{ArgType arg}}, \ldots) \\[4pt]
Class method:    & {\tt ClassName\_cls\_methodName\_}({\tt Class cls}, {\tt SEL \_cmd}, \\
                 & \quad {\meta{ArgType arg}}, \ldots) \\
\end{tabular}
\end{quote}

Examples:

\medskip
\noindent
\begin{tabular}{ll}
{\tt - (int)area}
  & becomes {\tt int Rectangle\_inst\_area\_(id self, SEL \_cmd)} \\[4pt]
{\tt - (id)initWithX:(int)x y:(int)y}
  & becomes {\tt id Rectangle\_inst\_initWithX\_y\_(id self, SEL \_cmd, int x, int y)}
  \\[4pt]
{\tt + (id)alloc}
  & becomes {\tt id Rectangle\_cls\_alloc\_(Class cls, SEL \_cmd)} \\
\end{tabular}

The first two parameters of every instance method implementation are
always {\tt self} (of type {\tt id}) and {\tt \_cmd} (of type {\tt SEL}).
Class methods receive {\tt cls} (of type {\tt Class}) and {\tt \_cmd}
instead.

% ========================================================================
\subsection{Method Definitions}
% ========================================================================

A method definition includes the full method body:
\begin{quote}
\begin{verbatim}
- (int)area {
    return width * height;
}
\end{verbatim}
\end{quote}

Within the method body, the preprocessor applies the following
transformations:

\begin{enumerate}
  \item The method is emitted as a C function with the naming
        convention described in Section~5.4.

  \item For instance methods, the implicit {\tt self} parameter is
        re-declared and cast to the class struct type:
        \begin{quote}
        {\tt struct Rectangle *self = (struct Rectangle *)\_obj;}
        \end{quote}
        This allows direct ivar access within the method body.

  \item Message expressions are translated to runtime calls
        (see Section~5.6).

  \item Bare ivar references of the form {\tt \_name} are rewritten to
        {\tt self->\_name} (see Section~5.7).

  \item {\tt @selector()} and {\tt @encode()} expressions are translated
        (see Sections~5.8 and~5.9).

  \item {\tt @autoreleasepool} blocks are expanded (see Section~5.10).
\end{enumerate}

% ========================================================================
\subsection{Message Expressions}
% ========================================================================

The primary mechanism for method invocation in Objective-C is the
message expression:
\begin{quote}
{\tt [}{\meta{receiver}} {\meta{selector}}: {\meta{argument}}{\tt ]}
\end{quote}

The preprocessor translates each message expression to a call to
{\tt objc\_msgSend()} or, for {\tt super} receivers,
{\tt objc\_msgSendSuper()}.

{\bf Input:}
\begin{quote}
{\tt [self draw];}
\end{quote}

{\bf Output:}
\begin{quote}
{\tt objc\_msgSend(self, sel\_registerName("draw"));}
\end{quote}

{\bf Input:}
\begin{quote}
{\tt [rect initWithX:10 y:20];}
\end{quote}

{\bf Output:}
\begin{quote}
{\tt objc\_msgSend(rect, sel\_registerName("initWithX:y:"), 10, 20);}
\end{quote}

Multi-keyword selectors are concatenated: the selector
{\tt initWithFoo:bar:} becomes the string {\tt ``initWithFoo:bar:''}.

% -----------------------------------------------------------------------
\subsubsection{Super Messages}
% -----------------------------------------------------------------------

When the receiver is the keyword {\tt super}, the message is dispatched
to the superclass implementation:

{\bf Input:}
\begin{quote}
{\tt [super draw];}
\end{quote}

{\bf Output:}
\begin{quote}
\begin{verbatim}
objc_msgSendSuper(self,
    class_getSuperclass(objc_getClass("Rectangle")),
    sel_registerName("draw"));
\end{verbatim}
\end{quote}

% -----------------------------------------------------------------------
\subsubsection{Class Messages}
% -----------------------------------------------------------------------

When the receiver is a class name (identified by a leading uppercase
letter), it is resolved at runtime:

{\bf Input:}
\begin{quote}
{\tt [Rectangle alloc];}
\end{quote}

{\bf Output:}
\begin{quote}
\begin{verbatim}
objc_msgSend((id)objc_getClass("Rectangle"),
    sel_registerName("alloc"));
\end{verbatim}
\end{quote}

% -----------------------------------------------------------------------
\subsubsection{Nested Expressions}
% -----------------------------------------------------------------------

Message expressions may be nested arbitrarily:

{\bf Input:}
\begin{quote}
{\tt [self initWithFrame:[NSScreen mainScreen]];}
\end{quote}

{\bf Output:}
\begin{quote}
\begin{verbatim}
objc_msgSend(self,
    sel_registerName("initWithFrame:"),
    objc_msgSend((id)objc_getClass("NSScreen"),
        sel_registerName("mainScreen")));
\end{verbatim}
\end{quote}

% -----------------------------------------------------------------------
\subsubsection{Heuristic Validation}
% -----------------------------------------------------------------------

Not every bracket expression in Objective-C source is a message
expression.  Array subscripts, for example, also use brackets:
\begin{quote}
{\tt int x = array[index];}
\end{quote}

The preprocessor applies a set of heuristic tests to distinguish message
expressions from C bracket expressions:

\begin{itemize}
  \item The receiver must begin with a letter or underscore.
  \item The selector must contain at least one alphabetic character.
  \item The receiver must not be a numeric literal.
\end{itemize}

When these tests fail, the bracket expression is passed through
unchanged to the C compiler.

% ========================================================================
\subsection{Implicit {\tt self} for Ivar Access}
% ========================================================================

Within instance method bodies, identifiers beginning with an underscore
that are not part of a larger expression are automatically rewritten to
reference {\tt self}.

{\bf Input:}
\begin{quote}
\begin{verbatim}
- (int)area {
    return _width * _height;
}
\end{verbatim}
\end{quote}

{\bf Output:}
\begin{quote}
\begin{verbatim}
int Rectangle_inst_area_(id self, SEL _cmd) {
    struct Rectangle *self = (struct Rectangle *)self;
    return self->_width * self->_height;
}
\end{verbatim}
\end{quote}

This rewrite applies only when the underscore-prefixed identifier
appears at the beginning of a new token---after whitespace, operators,
or the start of a line.  Identifiers such as {\tt my\_width} are not
affected.

The rewrite does not apply to class methods, which have no instance
context.

% ========================================================================
\subsection{Selector Expressions}
% ========================================================================

The {\tt @selector()} expression produces a selector value at compile
time:
\begin{quote}
{\tt @selector(}{\meta{methodName}}{\tt )}
\end{quote}

This is translated to a runtime call:
\begin{quote}
{\tt sel\_registerName(}{\meta{"methodName"}}{\tt )}
\end{quote}

Multi-word selectors are handled correctly:
\begin{quote}
\begin{tabular}{ll}
{\tt @selector(initWithFoo:bar:)}
  & becomes {\tt sel\_registerName("initWithFoo:bar:")} \\
\end{tabular}
\end{quote}

The selector name is trimmed of trailing whitespace.

% ========================================================================
\subsection{Encode Expressions}
% ========================================================================

The {\tt @encode()} expression yields a string containing the type
encoding of the given type:
\begin{quote}
\begin{tabular}{ll}
{\tt @encode(int)}
  & becomes {\tt "int"} \\[4pt]
{\tt @encode(NSString *)}
  & becomes {\tt "NSString *"} \\
\end{tabular}
\end{quote}

Nested parenthesized types are handled:
\begin{quote}
\begin{tabular}{ll}
{\tt @encode(struct \{ int x; \})}
  & becomes {\tt "struct \{ int x; \}"} \\
\end{tabular}
\end{quote}

% ========================================================================
\subsection{Autorelease Pools}
% ========================================================================

The {\tt @autoreleasepool} directive provides scoped memory management:
\begin{quote}
\begin{verbatim}
@autoreleasepool {
    id obj = [[NSObject alloc] init];
}
\end{verbatim}
\end{quote}

This is expanded to an explicit push/pop pair:
\begin{quote}
\begin{verbatim}
{ /* @autoreleasepool */
    void *_objc_apool = objc_autoreleasePoolPush();
    { /* @autoreleasepool */
        id obj = objc_msgSend(
            objc_msgSend((id)objc_getClass("NSObject"),
                sel_registerName("alloc")),
            sel_registerName("init"));
    }
    objc_autoreleasePoolPop(_objc_apool);
}
\end{verbatim}
\end{quote}

% ========================================================================
\subsection{Compatibility Aliases}
% ========================================================================

The {\tt @compatibility\_alias} directive introduces an alternative name
for an existing class:
\begin{quote}
{\tt @compatibility\_alias} {\meta{AliasName}} {\meta{OriginalName}};
\end{quote}

This is translated to a preprocessor macro:
\begin{quote}
{\tt \#define} {\meta{AliasName}} {\meta{OriginalName}} \hfill
  {\tt /* @compatibility\_alias */}
\end{quote}

% ========================================================================
\subsection{Access Control Keywords}
% ========================================================================

The following keywords control ivar and method visibility:

\medskip
\noindent
\begin{tabularx}{\textwidth}{lX}
{\tt @public}    & Ivars are accessible from any code. \\
{\tt @protected} & Ivars are accessible within the class and subclasses. \\
{\tt @private}   & Ivars are accessible only within the class. \\
{\tt @package}   & Ivars are accessible within the same framework. \\
\end{tabularx}

These keywords are converted to C comments in the preprocessor output.
Visibility enforcement is provided by the runtime.

Similarly, {\tt @optional} and {\tt @required} (used in protocol
declarations) are converted to comments.

% ========================================================================
\subsection{Exception Handling}
% ========================================================================

The preprocessor provides partial support for exception handling:
\begin{quote}
\begin{verbatim}
@try {
    ...
}
@catch (id exception) {
    ...
}
@finally {
    ...
}
\end{verbatim}
\end{quote}

The {\tt @throw} expression is translated to:
\begin{quote}
{\tt objc\_exception\_throw(}{\meta{expression}}{\tt )}
\end{quote}

The {\tt @try}, {\tt @catch}, and {\tt @finally} blocks are currently
emitted as C block structures with comment annotations.  Full exception
handling support requires the runtime's {\tt longjmp}-based exception
mechanism (see the \libobjcrt\ documentation).

% ========================================================================
\newpage
\section{C Preprocessor Directives}
% ========================================================================

% ========================================================================
\subsection{{\tt \#import} and {\tt \#include}}
% ========================================================================

The preprocessor provides special handling for {\tt \#import} and
{\tt \#include} directives with quoted filenames:
\begin{quote}
{\tt \#import "MyClass.h"}\\
{\tt \#include "MyClass.h"}
\end{quote}

When a quoted filename is given, the preprocessor opens the file and
recursively preprocesses its contents inline.  This provides the
semantics of {\tt \#import} (inclusion without re-inclusion) even for
files that do not contain their own include guards.

Include guard tracking prevents a file from being processed more than
once.  Up to 64 unique filenames may be tracked.

For angle-bracket includes, the standard behavior is preserved:
\begin{quote}
{\tt \#import <Foundation/Foundation.h>}
\end{quote}

This is passed through as a standard {\tt \#include} directive for the C
compiler to resolve using its own search path.

Include path search: when a quoted {\tt \#import} fails to open relative
to the current directory, the paths specified by {\tt -I} flags are
searched in order.

% ========================================================================
\subsection{The {\tt \#objc} Directive}
% ========================================================================

The {\tt \#objc} directive is a MinSTEP extension for including the
Objective-C standard library:
\begin{quote}
{\tt \#objc}
\end{quote}

When invoked without an argument, this includes the root header
{\tt objc/objc.h}, which declares the {\tt Object} base class and
standard Foundation types.

\begin{quote}
{\tt \#objc} {\meta{MyClass.h}}
\end{quote}

When invoked with a path argument, the specified file is recursively
preprocessed.

% ========================================================================
\subsection{Standard C Directives}
% ========================================================================

All standard C preprocessor directives are passed through to the
underlying C compiler without interpretation:

\medskip
\noindent
\begin{tabularx}{\textwidth}{lX}
{\tt \#define}   & Define a macro. \\
{\tt \#undef}    & Undefine a macro. \\
{\tt \#if}       & Begin conditional compilation. \\
{\tt \#ifdef}    & Test whether a macro is defined. \\
{\tt \#ifndef}   & Test whether a macro is not defined. \\
{\tt \#elif}     & Else-if conditional branch. \\
{\tt \#else}     & Else conditional branch. \\
{\tt \#endif}    & End conditional compilation. \\
{\tt \#pragma}   & Compiler-specific directives. \\
{\tt \#error}    & Emit a compilation error. \\
{\tt \#warning}  & Emit a compilation warning. \\
{\tt \#line}     & Override line number information. \\
\end{tabularx}

The preprocessor does not expand macros, evaluate constant expressions,
or perform token pasting.  These tasks are left to the C compiler.

% ========================================================================
\newpage
\section{The Runtime Library}
% ========================================================================

% ========================================================================
\subsection{Overview}
% ========================================================================

The MinSTEP runtime library, \libobjcrt, provides the runtime support
required by Objective-C programs.  It is written in pure C and linkable
with any ANSI C compiler.

The runtime provides:

\begin{itemize}
  \item Dynamic class registration and lookup
  \item Message dispatch ({\tt objc\_msgSend})
  \item Selector registration and lookup
  \item Protocol management
  \item Manual reference counting ({\tt retain}/{\tt release}/
        {\tt autorelease})
  \item Autorelease pool management
  \item Exception handling primitives
  \item Property and instance variable metadata
\end{itemize}

The runtime is initialized automatically at load time via a constructor
function; no explicit initialization call is required.

% ========================================================================
\subsection{Message Dispatch}
% ========================================================================

All Objective-C method calls are routed through the dispatch function:
\begin{quote}
{\tt id objc\_msgSend(id self, SEL \_cmd, ...)}
\end{quote}

The runtime locates the method implementation for the given receiver and
selector by walking the class hierarchy.  If the method is not found, the
{\tt forwardInvocation:} handler is invoked.

For super messages, a separate function is provided:
\begin{quote}
\begin{tabular}{l}
{\tt id objc\_msgSendSuper(id self, struct objc\_class *superclass,}\\
\quad {\tt SEL \_cmd, ...)}
\end{tabular}
\end{quote}

This dispatches the message starting from the given superclass,
bypassing the receiver's own class.

% ========================================================================
\subsection{Class and Selector Registry}
% ========================================================================

Classes are registered at load time via
{\tt objc\_allocateClassPair()} or the {\tt @implementation} directive.
The runtime maintains a hash table of all registered classes, indexed by
name.

Selectors are interned strings.  The {\tt sel\_registerName()} function
returns the unique {\tt SEL} value for a given selector name string.
Multiple calls with the same name return the same pointer value, allowing
pointer comparison for selector equality.

% ========================================================================
\subsection{Memory Management}
% ========================================================================

The MinSTEP runtime implements manual reference counting (MRC).  Each
object's reference count is tracked in an internal hash table.

\medskip
\noindent
\begin{tabularx}{\textwidth}{lX}
{\tt id objc\_retain(id obj)}
  & Increment the reference count. \\[4pt]
{\tt id objc\_release(id obj)}
  & Decrement the reference count; deallocate at zero. \\[4pt]
{\tt id objc\_autorelease(id obj)}
  & Register {\tt obj} for release at the end of the current autorelease
    pool scope. \\
\end{tabularx}

Autorelease pools are managed with push/pop pairs:
\begin{quote}
\begin{tabular}{l}
{\tt void *objc\_autoreleasePoolPush(void)}\\
{\tt void\ \ objc\_autoreleasePoolPop(void *pool)}
\end{tabular}
\end{quote}

When compiled with ARC (Automatic Reference Counting), the
{\tt retain}/{\tt release}/{\tt autorelease} calls are omitted; the
compiler inserts them automatically.  The runtime detects ARC
compilation via the {\tt \_\_has\_feature(objc\_arc)} preprocessor test.

% ========================================================================
\subsection{Instance Variables}
% ========================================================================

Instance variables are laid out in the class struct immediately after
the {\tt isa} and {\tt \_super} pointers.  The runtime provides
functions for accessing ivars by name:

\medskip
\noindent
\begin{tabularx}{\textwidth}{lX}
{\tt Ivar class\_getInstanceVariable(Class cls, const char *name)}
  & Look up an ivar by name. \\[4pt]
{\tt id object\_getIvar(id obj, Ivar ivar)}
  & Read an ivar value. \\[4pt]
{\tt void object\_setIvar(id obj, Ivar ivar, id value)}
  & Write an ivar value. \\
\end{tabularx}

The preprocessor's {\tt self->\_name} rewriting (Section~5.7) provides
direct struct access, which is more efficient than the runtime ivars API
for compile-time-known ivars.

% ========================================================================
\subsection{Protocols}
% ========================================================================

Protocols are registered with the runtime and may be queried at runtime:

\medskip
\noindent
\begin{tabularx}{\textwidth}{lX}
{\tt Protocol *objc\_allocateProtocol(const char *name)}
  & Allocate a new protocol. \\[4pt]
{\tt void objc\_registerProtocol(Protocol *proto)}
  & Register an allocated protocol. \\[4pt]
{\tt Protocol *objc\_getProtocol(const char *name)}
  & Look up a registered protocol by name. \\
\end{tabularx}

Classes adopt protocols via the runtime's class registration API.  The
preprocessor emits {\tt Protocol} pointer variables for each
{\tt @protocol} declaration; the runtime initializes these at load time.

% ========================================================================
\newpage
\section{Building the Toolchain}
% ========================================================================

% ========================================================================
\subsection{Prerequisites}
% ========================================================================

Building the MinSTEP Objective-C toolchain requires:

\begin{itemize}
  \item A C99-conforming compiler ({\tt gcc} or {\tt clang})
  \item GNU make
  \item Standard POSIX utilities
\end{itemize}

No external libraries are required beyond the C standard library and the
POSIX runtime.

% ========================================================================
\subsection{Build Commands}
% ========================================================================

To build the entire toolchain:
\begin{quote}
{\tt \$ make}
\end{quote}

This produces:
\begin{itemize}
  \item {\tt objcc} \quad The preprocessor executable.
  \item {\tt build/libobjcrt.a} \quad The runtime static library.
\end{itemize}

To build individual components:
\begin{quote}
{\tt \$ make objcc} \quad\quad Build only the preprocessor.\\
{\tt \$ make libobjcrt} \quad Build only the runtime library.
\end{quote}

To install:
\begin{quote}
{\tt \$ make install PREFIX=/usr/local}
\end{quote}

This installs:
\begin{quote}
\begin{tabular}{ll}
{\tt /usr/local/bin/objcc} & The preprocessor. \\
{\tt /usr/local/lib/libobjcrt.a} & Runtime static library. \\
{\tt /usr/local/lib/libobjcrt.o} & Runtime object file. \\
{\tt /usr/local/include/objc/objc.h} & Standard library header. \\
{\tt /usr/local/include/objc/objcrt.h} & Runtime API header. \\
\end{tabular}
\end{quote}

% ========================================================================
\subsection{Linking}
% ========================================================================

To link an Objective-C program:
\begin{quote}
{\tt \$ objcc -o myapp main.m -lobjcrt}
\end{quote}

If the runtime is installed in a non-standard location:
\begin{quote}
{\tt \$ objcc -L /opt/minstep/lib -l objcrt -o myapp main.m}
\end{quote}

% ========================================================================
\newpage
\section{Error Handling}
% ========================================================================

% ========================================================================
\subsection{Diagnostics}
% ========================================================================

The preprocessor emits error messages to the standard error output when
it encounters conditions it cannot handle:

\begin{itemize}
  \item Unterminated string literals or comments.
  \item Unbalanced braces within method definitions.
  \item {\tt @end} encountered without a matching {\tt @interface} or
        {\tt @implementation}.
  \item File not found for {\tt \#import} or {\tt \#include} directives.
\end{itemize}

Error messages include the input filename and line number:
\begin{quote}
{\tt MyClass.m:42: error: unterminated string literal}
\end{quote}

The preprocessor continues processing after errors when possible,
emitting as much valid output as it can.

% ========================================================================
\newpage
\section{Implementation Notes}
% ========================================================================

% ========================================================================
\subsection{Translation Strategy}
% ========================================================================

The preprocessor operates as a character-level stream processor.  It
does not construct an abstract syntax tree or token stream.  Instead, it
reads characters one at a time and dispatches on trigger characters:

\medskip
\noindent
\begin{tabular}{ll}
{\tt \#} & $\rightarrow$ C preprocessor directive \\
{\tt @} & $\rightarrow$ Objective-C keyword or expression \\
{\tt [} & $\rightarrow$ Possibly a message expression \\
{\tt +} & $\rightarrow$ Possibly a method definition or declaration \\
{\tt -} & $\rightarrow$ Possibly a method definition or declaration \\
\end{tabular}

This approach is simple and efficient.  The trade-off is that the
preprocessor cannot perform sophisticated semantic analysis.  It relies
on heuristic tests to distinguish Objective-C constructs from C
constructs that use similar syntax (particularly bracket expressions).

% ========================================================================
\subsection{Limitations}
% ========================================================================

The following Objective-C features are not currently supported by the
preprocessor:

\begin{itemize}
  \item {\tt for\ldots in} fast enumeration.
  \item Subscript syntax ({\tt obj[key]} and {\tt obj[key] = value}).
  \item Blocks (closures).
  \item Typed {\tt @encode()} for all types.
  \item Full {\tt @try/@catch/@finally} exception handling (partially
        stubbed).
  \item {\tt @protocol()} runtime protocol expressions.
  \item Class extensions in {\tt @interface} (informal protocols with
        category syntax are supported).
  \item {\tt @required} and {\tt @optional} method modifiers (emitted
        as comments only).
\end{itemize}

The runtime library does not yet implement:

\begin{itemize}
  \item Method cache lookup (methods are found by linear search of the
        superclass chain).
  \item Weak references.
  \item Associated objects.
  \item Tagged pointers.
  \item Non-fragile ABI.
\end{itemize}

% ========================================================================
\subsection{Performance}
% ========================================================================

The preprocessor is designed for correctness and simplicity, not for
speed.  It processes input in a single pass with no lookahead beyond one
character.  The generated C code is functionally equivalent to the
Objective-C input but may differ in structure.

The runtime's message dispatch walks the class hierarchy linearly to
locate method implementations.  For programs that make frequently-invoked
method calls, this may be slower than a cached dispatch.  A method cache
is planned for a future release.

% ========================================================================
\newpage
\section*{Appendix~A:\quad Translation Reference}
\addcontentsline{toc}{section}{Appendix A: Translation Reference}
% ========================================================================

The following table summarizes the translation of each Objective-C
construct.

\bigskip
\noindent
\begin{longtable}{>{\tt}l >{\tt}l}
\hline
\multicolumn{1}{l}{\bf Input Construct} &
\multicolumn{1}{l}{\bf Output C} \\
\hline
\endfirsthead
\hline
\multicolumn{1}{l}{\bf Input Construct} &
\multicolumn{1}{l}{\bf Output C} \\
\hline
\endhead
\hline
\endfoot
@
interface Foo : Bar \{ \} & struct Foo \{ \ldots \}; \\
@implementation Foo & /\ast\ @implementation\ Foo\ \ast/ \\
@protocol MyProto & static Protocol\ *\_p = NULL; \\
@property (nonatomic) int x; & --(int)x; --(void)setX:(int)v; \\
@synthesize x = \_x; & getter/setter bodies \\
@dynamic x; & /\ast\ comment\ \ast/ \\
@class Foo; & struct Foo; \\
@compatibility\_alias A B; & \#define A B \\
@selector(foo) & sel\_registerName("foo") \\
@encode(int) & "int" \\
@autoreleasepool \{ \} & push/pop block \\
[self foo:bar] & objc\_msgSend(self, sel("foo:"), bar) \\
super foo] & objc\_msgSendSuper(\ldots) \\
\_ivar & self-->\_ivar \\
+ (id)alloc & ClassName\_cls\_alloc\_(\ldots) \\
--(int)getX & ClassName\_inst\_getX\_(\ldots) \\
\#import "file.h" & recursive inline inclusion \\
\#objc & \#include "objc/objc.h" \\
\end{longtable}

% ========================================================================
\section*{Appendix~B:\quad Configuration Limits}
\addcontentsline{toc}{section}{Appendix B: Configuration Limits}
% ========================================================================

The following limits are compile-time constants in the preprocessor.

\bigskip
\noindent
\begin{tabularx}{\textwidth}{lrlX}
\hline
\bf Symbol & \bf Value & & \bf Description \\
\hline
{\tt MAX\_TOKEN}       & 4{,}096 & & Maximum token length. \\
{\tt MAX\_LINE}        & 8{,}192 & & Maximum line length. \\
{\tt MAX\_CLASSES}     & 512     & & Maximum tracked class names. \\
{\tt MAX\_PROTOCOLS}   & 256     & & Maximum tracked protocol names. \\
{\tt MAX\_PROPERTIES}  & 1{,}024 & & Maximum tracked properties. \\
{\tt MAX\_INCLUDES}    & 64      & & Maximum include guard entries. \\
{\tt MAX\_INCLUDE\_PATHS} & 64  & & Maximum {\tt -I} search paths. \\
\hline
\end{tabularx}

\vfill

% -----------------------------------------------------------------------
%  Colophon
% -----------------------------------------------------------------------

\begin{center}
\small\itshape
This document was typeset in \TeX.\\
The body text is set in Computer Modern Roman,\\
and code examples in Computer Modern Typewriter.
\end{center}

\end{document}
